\documentclass[10pt,a4paper]{amsart}
\usepackage{amsmath,amssymb,amsthm,graphicx}
\usepackage{verbatim}
\newenvironment{reference}{\comment}{\endcomment}
%\newenvironment{reference}{}{}
\newenvironment{slogan}{\comment}{\endcomment}
\newenvironment{history}{\comment}{\endcomment}

% For commutative diagrams we use Xy-pic
\usepackage[all]{xy}

% We use 2cell for 2-commutative diagrams.
\xyoption{2cell}
% Theorem environments.
%
\theoremstyle{plain}
\newtheorem{theorem}[subsection]{Theorem}
\newtheorem{proposition}[subsection]{Proposition}
\newtheorem{lemma}[subsection]{Lemma}

\theoremstyle{definition}
\newtheorem{definition}[subsection]{Definition}
\newtheorem{example}[subsection]{Example}
\newtheorem{exercise}[subsection]{Exercise}
\newtheorem{situation}[subsection]{Situation}

\theoremstyle{remark}
\newtheorem{remark}[subsection]{Remark}
\newtheorem{remarks}[subsection]{Remarks}

\numberwithin{equation}{subsection}

% Macros
%
\def\lim{\mathop{\mathrm{lim}}\nolimits}
\def\colim{\mathop{\mathrm{colim}}\nolimits}
\def\Spec{\mathop{\mathrm{Spec}}}
\def\Hom{\mathop{\mathrm{Hom}}\nolimits}
\def\Ext{\mathop{\mathrm{Ext}}\nolimits}
\def\SheafHom{\mathop{\mathcal{H}\!\mathit{om}}\nolimits}
\def\SheafExt{\mathop{\mathcal{E}\!\mathit{xt}}\nolimits}
\def\Sch{\mathit{Sch}}
\def\Mor{\mathop{\mathrm{Mor}}\nolimits}
\def\Ob{\mathop{\mathrm{Ob}}\nolimits}
\def\Sh{\mathop{\mathit{Sh}}\nolimits}
\def\NL{\mathop{N\!L}\nolimits}
\def\CH{\mathop{\mathrm{CH}}\nolimits}
\def\proetale{{pro\text{-}\acute{e}tale}}
\def\etale{{\acute{e}tale}}
\def\QCoh{\mathit{QCoh}}
\def\Ker{\mathop{\mathrm{Ker}}}
\def\Im{\mathop{\mathrm{Im}}}
\def\Coker{\mathop{\mathrm{Coker}}}
\def\Coim{\mathop{\mathrm{Coim}}}

% Boxtimes
%
\DeclareMathSymbol{\boxtimes}{\mathbin}{AMSa}{"02}

%
% Macros for moduli stacks/spaces
%
\def\QCohstack{\mathcal{QC}\!\mathit{oh}}
\def\Cohstack{\mathcal{C}\!\mathit{oh}}
\def\Spacesstack{\mathcal{S}\!\mathit{paces}}
\def\Quotfunctor{\mathrm{Quot}}
\def\Hilbfunctor{\mathrm{Hilb}}
\def\Curvesstack{\mathcal{C}\!\mathit{urves}}
\def\Polarizedstack{\mathcal{P}\!\mathit{olarized}}
\def\Complexesstack{\mathcal{C}\!\mathit{omplexes}}
% \Pic is the operator that assigns to X its picard group, usage \Pic(X)
% \Picardstack_{X/B} denotes the Picard stack of X over B
% \Picardfunctor_{X/B} denotes the Picard functor of X over B
\def\Pic{\mathop{\mathrm{Pic}}\nolimits}
\def\Picardstack{\mathcal{P}\!\mathit{ic}}
\def\Picardfunctor{\mathrm{Pic}}
\def\Deformationcategory{\mathcal{D}\!\mathit{ef}}
\usepackage{cleveref}
\crefname{originallemma}{original lemma}{original lemmas}
\Crefname{originallemma}{Original lemma}{Original lemmas}
\makeatletter
\newlabel{algebra-lemma-etale}{{00U2}{source}{}{sourcepointer-1}{}}
\newlabel{algebra-lemma-etale@cref}{{[originallemma][0][]00U2}{[1][0][]source}{}{}{}}
\makeatother
\crefname{originallemma}{original lemma}{original lemmas}
\Crefname{originallemma}{Original lemma}{Original lemmas}
\makeatletter
\newlabel{algebra-lemma-lift-etale}{{04D1}{source}{}{sourcepointer-2}{}}
\newlabel{algebra-lemma-lift-etale@cref}{{[originallemma][0][]04D1}{[1][0][]source}{}{}{}}
\makeatother
\crefname{originallemma}{original lemma}{original lemmas}
\Crefname{originallemma}{Original lemma}{Original lemmas}
\makeatletter
\newlabel{algebra-lemma-map-between-etale}{{00U7}{source}{}{sourcepointer-3}{}}
\newlabel{algebra-lemma-map-between-etale@cref}{{[originallemma][0][]00U7}{[1][0][]source}{}{}{}}
\makeatother
\crefname{originallemma}{original lemma}{original lemmas}
\Crefname{originallemma}{Original lemma}{Original lemmas}
\makeatletter
\newlabel{algebra-lemma-pure}{{04PS}{source}{}{sourcepointer-4}{}}
\newlabel{algebra-lemma-pure@cref}{{[originallemma][0][]04PS}{[1][0][]source}{}{}{}}
\makeatother
\crefname{originaldefinition}{original definition}{original definitions}
\Crefname{originaldefinition}{Original definition}{Original definitions}
\makeatletter
\newlabel{categories-definition-directed}{{002V}{source}{}{sourcepointer-5}{}}
\newlabel{categories-definition-directed@cref}{{[originaldefinition][0][]002V}{[1][0][]source}{}{}{}}
\makeatother
\crefname{originallemma}{original lemma}{original lemmas}
\Crefname{originallemma}{Original lemma}{Original lemmas}
\makeatletter
\newlabel{categories-lemma-adjoint-fully-faithful}{{07RB}{source}{}{sourcepointer-6}{}}
\newlabel{categories-lemma-adjoint-fully-faithful@cref}{{[originallemma][0][]07RB}{[1][0][]source}{}{}{}}
\makeatother
\crefname{originalreference}{original reference}{original references}
\Crefname{originalreference}{Original reference}{Original references}
\makeatletter
\newlabel{lemma-ind-etale-permanence}{{097M}{source}{}{sourcepointer-7}{}}
\newlabel{lemma-ind-etale-permanence@cref}{{[originalreference][0][]097M}{[1][0][]source}{}{}{}}
\makeatother
\title{Base change along an ideal on ind-etale algebras has a fully faithful right adjoint}
\author{The Stacks Project Authors}
\date{stacks-project:proetale.tex:local-snapshot-20261012; GNU FDL 1.2 or later}
\begin{document}
\maketitle
\noindent\textit{Editorial note.} All delivered mathematical text is the author's own, copied from the pinned Stacks Project chapter file proetale.tex; no statement or proof step was rewritten, reordered or removed.\par\medskip
\noindent\textit{Editorial note.} The counted proof is the author's complete argument: the directedness of the category of factorizations through etale algebras, the identification of the colimit modulo the ideal with the given algebra via flatness and a localization argument, the computation of the Hom sets establishing the adjunction, and the closing full faithfulness.\par\medskip
\noindent\textit{Editorial note.} The label lemma-lift-ind-etale names the counted statement, and the proof environment immediately following it in the source is delivered verbatim as its proof.\par\medskip
\noindent\textit{Editorial note.} The counted statement is carried by the author's own lemma environment in the chapter source.\par\medskip
\noindent\textit{Editorial note.} Pre-open gate: python operations/agent-authoring/tools/precard\_gate.py --source proetale.tex --label lemma-lift-ind-etale --start 1437 --end 1516 --local-id 3577 returned RESULT=OK (L3577 free, occupancy FREE, public interval clear). Candidate taken from sources/stacks-proetale/candidates-halfA.json (new exclusive pool assigned by the Lead on 2026-10-13).\par\medskip
\noindent\textit{Editorial note.} The proof contains the parenthetical '(We ignore some set theoretical issues in this proof.)'; this is a standard set-theoretic caveat about the size of the category of factorizations and does not omit any step of the mathematical argument, so the statement is counted with the parenthetical recorded here verbatim.\par\medskip
\noindent\textit{Editorial note.} The statement and proof are copied from the pinned chapter file; the cited inputs are kept as external pointers with their canonical Stacks tags.\par\medskip
\noindent\textit{Allowed standard prerequisites (editorial).}\begin{itemize}
\item Characterization of directed categories. Categories, the definition of a directed category. Supplies the three properties (nonempty, upper bounds for pairs, coequalizers) checked for the category of factorizations.
\item Stability of etale algebras under tensor products, quotients by coequalizers and localization. Algebra, the lemmas that etale algebras are stable under tensor product, that a map between etale algebras is etale, and the lifting of etale algebras along a surjection of rings, together with the lemma on pure ring maps used to identify a localization. Supply the objects of the category of factorizations, the surjectivity of C -> C-bar, the identification of C/IC with C-bar and the lifting of the etale pieces.
\item Permanence of the ind-etale property. Same chapter, the lemma that the ind-etale property is preserved under the operations used in the proof. Supplies the flatness of C/IC -> C-bar, which together with the pureness lemma gives the localization statement.
\item Full faithfulness of a functor with an adjoint whose counit is the identity. Categories, the lemma that a functor with a left adjoint such that the unit is an isomorphism is fully faithful. Supplies the conclusion that the right adjoint v is fully faithful because the composite of the base change functor with v is the identity.
\end{itemize}\medskip
\ifdefined\hypertarget\hypertarget{sourcepointer-1}{}\fi
\noindent\textit{Source reference.} Original lemma 00U2: sources/stacks-chow/tags.txt, source lines 1096--1096. Prerequisite pointer invoked at proof line(s) 1467; it names a standard fact of a sibling chapter that the counted argument uses, and every step of the counted proof built on it is retained verbatim in the delivered body.\par\medskip
\ifdefined\hypertarget\hypertarget{sourcepointer-2}{}\fi
\noindent\textit{Source reference.} Original lemma 04D1: sources/stacks-chow/tags.txt, source lines 5711--5711. Prerequisite pointer invoked at proof line(s) 1481; it names a standard fact of a sibling chapter that the counted argument uses, and every step of the counted proof built on it is retained verbatim in the delivered body.\par\medskip
\ifdefined\hypertarget\hypertarget{sourcepointer-3}{}\fi
\noindent\textit{Source reference.} Original lemma 00U7: sources/stacks-chow/tags.txt, source lines 1101--1101. Prerequisite pointer invoked at proof line(s) 1474; it names a standard fact of a sibling chapter that the counted argument uses, and every step of the counted proof built on it is retained verbatim in the delivered body.\par\medskip
\ifdefined\hypertarget\hypertarget{sourcepointer-4}{}\fi
\noindent\textit{Source reference.} Original lemma 04PS: sources/stacks-chow/tags.txt, source lines 6125--6125. Prerequisite pointer invoked at proof line(s) 1487; it names a standard fact of a sibling chapter that the counted argument uses, and every step of the counted proof built on it is retained verbatim in the delivered body.\par\medskip
\ifdefined\hypertarget\hypertarget{sourcepointer-5}{}\fi
\noindent\textit{Source reference.} Original definition 002V: sources/stacks-chow/tags.txt, source lines 112--112. Prerequisite pointer invoked at proof line(s) 1463; it names a standard fact of a sibling chapter that the counted argument uses, and every step of the counted proof built on it is retained verbatim in the delivered body.\par\medskip
\ifdefined\hypertarget\hypertarget{sourcepointer-6}{}\fi
\noindent\textit{Source reference.} Original lemma 07RB: sources/stacks-chow/tags.txt, source lines 9861--9861. Prerequisite pointer invoked at proof line(s) 1515; it names a standard fact of a sibling chapter that the counted argument uses, and every step of the counted proof built on it is retained verbatim in the delivered body.\par\medskip
\ifdefined\hypertarget\hypertarget{sourcepointer-7}{}\fi
\noindent\textit{Source reference.} Original reference 097M: sources/stacks-proetale/pinned-original-source.tex, source lines 1417--1417. Prerequisite pointer invoked at proof line(s) 1484; it names a standard fact of the same chapter that the counted argument uses, and every step of the counted proof built on it is retained verbatim in the delivered body.\par\medskip
\begin{lemma}
\label{lemma-lift-ind-etale}
Let $A$ be a ring and let $I \subset A$ be an ideal. The base change functor
$$
\text{ind-\'etale }A\text{-algebras}
\longrightarrow
\text{ind-\'etale }A/I\text{-algebras},\quad
C \longmapsto C/IC
$$
has a fully faithful right adjoint $v$. In particular, given
an ind-\'etale $A/I$-algebra $\overline{C}$ there exists
an ind-\'etale $A$-algebra $C = v(\overline{C})$ such that
$\overline{C} = C/IC$.
\end{lemma}

\begin{proof}
Let $\overline{C}$ be an ind-\'etale $A/I$-algebra.
Consider the category $\mathcal{C}$ of factorizations
$A \to B \to \overline{C}$ where $A \to B$ is \'etale.
(We ignore some set theoretical issues in this proof.)
We will show that this category is directed and that
$C = \colim_\mathcal{C} B$ is an ind-\'etale $A$-algebra
such that $\overline{C} = C/IC$.

\medskip\noindent
We first prove that $\mathcal{C}$ is directed
(Categories, Definition \ref{categories-definition-directed}).
The category is nonempty as $A \to A \to \overline{C}$ is an object.
Suppose that $A \to B \to \overline{C}$ and $A \to B' \to \overline{C}$
are two objects of $\mathcal{C}$. Then $A \to B \otimes_A B' \to \overline{C}$
is another (use Algebra, Lemma \ref{algebra-lemma-etale}).
Suppose that $f, g : B \to B'$ are two maps between
objects $A \to B \to \overline{C}$ and $A \to B' \to \overline{C}$
of $\mathcal{C}$. Then a coequalizer is
$A \to B' \otimes_{f, B, g} B' \to \overline{C}$.
This is an object of $\mathcal{C}$ by
Algebra, Lemmas \ref{algebra-lemma-etale} and
\ref{algebra-lemma-map-between-etale}.
Thus the category $\mathcal{C}$ is directed.

\medskip\noindent
Write $\overline{C} = \colim \overline{B_i}$ as a filtered colimit with
$\overline{B_i}$ \'etale over $A/I$. For every $i$
there exists $A \to B_i$ \'etale with $\overline{B_i} = B_i/IB_i$, see
Algebra, Lemma \ref{algebra-lemma-lift-etale}.
Thus $C \to \overline{C}$ is surjective.
Since $C/IC \to \overline{C}$ is ind-\'etale
(Lemma \ref{lemma-ind-etale-permanence})
we see that it is flat. Hence $\overline{C}$ is a localization of
$C/IC$ at some multiplicative subset $S \subset C/IC$
(Algebra, Lemma \ref{algebra-lemma-pure}).
Take an $f \in C$ mapping to an element of $S \subset C/IC$.
Choose $A \to B \to \overline{C}$ in $\mathcal{C}$ and $g \in B$
mapping to $f$ in the colimit. Then we see that $A \to B_g \to \overline{C}$
is an object of $\mathcal{C}$ as well. Thus $f$ is an invertible
element of $C$. It follows that $C/IC = \overline{C}$.

\medskip\noindent
Next, we claim that for an ind-\'etale algebra $D$ over $A$ we have
$$
\Mor_A(D, C) = \Mor_{A/I}(D/ID, \overline{C})
$$
Namely, let $D/ID \to \overline{C}$ be an $A/I$-algebra map.
Write $D = \colim_{i \in I} D_i$ as a colimit over a directed set $I$
with $D_i$ \'etale over $A$. By choice of $\mathcal{C}$
we obtain a transformation $I \to \mathcal{C}$ and hence a map
$D \to C$ compatible with maps to $\overline{C}$. Whence the claim.

\medskip\noindent
It follows that the functor $v$ defined by the rule
$$
\overline{C} 
\longmapsto
v(\overline{C}) = \colim_{A \to B \to \overline{C}} B
$$
is a right adjoint to the base change functor $u$ as required by the lemma.
The functor $v$ is fully faithful because
$u \circ v = \text{id}$ by construction, see
Categories, Lemma \ref{categories-lemma-adjoint-fully-faithful}.
\end{proof}

\end{document}
